\subsection{Proof by Division into Cases}
\label{subsections:techniques:universal:cases}



\paragraph{Description.}
A proof by division into cases proves a mathematical statement by splitting %
it into smaller statements and proving each one. %
Precisely, it proves a universal statement, $\forall s \in S,\ P(s)$, %
by dividing $S$ into subsets $S_1, \dots, S_n$ such that %
$S = S_1 \cup \dots \cup S_n$ and showing that $P(s)$ is true for each %
$s \in S_i$. %
Generally, a proof by division into cases mirrors the following format:
\begin{enumerate}{
    \item \textbf{Setup.} %
        Let $s \in S$ be arbitrary and divide $S$ into %
        $S_1, \dots, S_n$ such that $S = S_1 \cup \dots \cup S_n$.
    \item \textbf{Verification.} %
        For each $S_i$, suppose $s \in S_i$ and show that $P(s)$ is true.
    \item \textbf{Conclusion.} %
        Conclude the proof, optionally by summarizing what has been found.
}
\end{enumerate}



\paragraph{Advice.}
\begin{itemize}{
    \item Use this proof when no single argument works for every $s \in S$ %
        but one works within each $S_i$.
    \item Cover every element of $S$, as any $s$ in no case leaves the %
        statement unproven.
    \item To find cases, identify a property that divides $S$ into %
        subsets, such as parity or sign.
}
\end{itemize}



\paragraph{Example.}
\begin{exercise}{
    Prove that $\forall z \in \Z,\ \exists m \in \Z$ such that %
    $z^2 + z = 2m$.
}
\end{exercise}
\begin{proof}{
    This statement can be proven by division into cases.

    \begin{annotatedpanel}[Setup.]{
        Let $z \in \Z$. %
        Then $z$ is either even or odd.
    }{
        \textbf{Step 1.} Let $s \in S$ be arbitrary and divide $S$ into %
        cases. %
        Here, $S = \Z$, $s = z$, $S_1$ is the even integers, %
        and $S_2$ is the odd integers.
    }
    \end{annotatedpanel}

    \begin{annotatedpanel}[Case 1.]{
        Suppose $z$ is even. %
        Then $z = 2k$ with $k \in \Z$. %
        Thus, %
        \begin{align*}{
            z^2 + z &= (2k)^2 + 2k\\
                    &= 4k^2 + 2k\\
                    &= 2(2k^2 + k).
        }
        \end{align*}
        Let $m \in \Z$ such that $m = 2k^2 + k$. %
        Then $z^2 + z = 2(2k^2 + k) = 2m$.
    }{
        \textbf{Step 2.} Suppose $s \in S_1$ and show that $P(s)$ is true. %
        Here, $P(z)$ is the statement $\exists m \in \Z$ such that %
        \mbox{$z^2 + z = 2m$}.
    }
    \end{annotatedpanel}

    \begin{annotatedpanel}[Case 2.]{
        Suppose $z$ is odd. %
        Then $z = 2k + 1$ with $k \in \Z$. %
        Thus, %
        \begin{align*}{
            z^2 + z &= (2k + 1)^2 + (2k + 1)\\
                    &= 4k^2 + 4k + 1 + 2k + 1\\
                    &= 4k^2 + 6k + 2\\
                    &= 2(2k^2 + 3k + 1).
        }
        \end{align*}
        Let $m \in \Z$ such that $m = 2k^2 + 3k + 1$. %
        Then $z^2 + z = 2(2k^2 + 3k + 1) = 2m$.
    }{
        \textbf{Step 2.} Repeat for the remaining case. %
        Here, $s \in S_2$.
    }
    \end{annotatedpanel}

    \begin{annotatedpanel}[Conclusion.]{
        In either case, $z^2 + z = 2m$. %
        Thus, $\exists m \in \Z$ such that \mbox{$z^2 + z = 2m$}. %
        Therefore, the statement holds by division into cases. \qedhere
    }{
        \textbf{Step 3.} Conclude the proof. %
        Here, the summary states what has been found: %
        $P(z)$ holds in both cases.
    }
    \end{annotatedpanel}
}
\end{proof}